\documentclass[reqno]{amsart}
\usepackage[T1]{fontenc}
\usepackage{amsmath,amssymb,amsthm}
\usepackage{xr-hyper}
\usepackage[hidelinks]{hyperref}
\externaldocument[sm-]{20261125T000000--slope-method__workbook}
\externaldocument[ds-]{20261125T004000--denominators-and-sha__workbook}
\externaldocument[ts-]{20261125T004100--tate-twists-and-sha__workbook}
\externaldocument[cg-]{20261125T014100--congruence-modules__workbook}
\externaldocument[ei-]{20261125T014500--eisenstein-congruences-and-selmer-classes__workbook}
\externaldocument[pc-]{20261125T014700--pseudocharacters-and-the-length-bound__workbook}
\externaldocument[gb-]{20261125T011000--galois-representations-black-box__workbook}

\newcounter{chap}
% Chapter numbers, generated from 20261125T012000--grh-workbook__guide.org
% by scripts/make-chapters.py.  Do not edit: change the order of the
% index and rerun the script (./preview-site.sh and CI do this).
%
% Usage in a sheet, after \newcounter{chap}:
%   % Chapter numbers, generated from 20261125T012000--grh-workbook__guide.org
% by scripts/make-chapters.py.  Do not edit: change the order of the
% index and rerun the script (./preview-site.sh and CI do this).
%
% Usage in a sheet, after \newcounter{chap}:
%   % Chapter numbers, generated from 20261125T012000--grh-workbook__guide.org
% by scripts/make-chapters.py.  Do not edit: change the order of the
% index and rerun the script (./preview-site.sh and CI do this).
%
% Usage in a sheet, after \newcounter{chap}:
%   \input{chapters}\chapnumber{<sheet ID>}
\makeatletter
\providecommand\chapters@set[2]{\@namedef{chapters@#1}{#2}}
\chapters@set{20261125T001000}{1}% zeros-to-ramp-sums
\chapters@set{20261125T001100}{2}% inert-primes-and-estimate-b
\chapters@set{20261125T005000}{3}% galois-fourier-and-averaging
\chapters@set{20261125T005500}{4}% columns-from-a-sphere
\chapters@set{20261125T000000}{5}% slope-method
\chapters@set{20261126T030100}{6}% a-twist-for-each-prime
\chapters@set{20261125T003000}{7}% the-chain-and-the-square-root
\chapters@set{20261125T010000}{8}% rank-lemma-and-zero-estimate
\chapters@set{20261125T010500}{9}% poincare-series-and-near-orthogonality
\chapters@set{20261124T223000}{10}% elliptic-curves-for-the-proof
\chapters@set{20261125T009500}{11}% fock-model-and-theta-intertwiners
\chapters@set{20261124T223100}{12}% complex-multiplication-and-inert-reduction
\chapters@set{20261125T008000}{13}% one-cm-determinant
\chapters@set{20261124T214500}{14}% supersingular-theta-divisibility
\chapters@set{20261124T214600}{15}% the-isotropic-plane
\chapters@set{20261125T013700}{16}% pairings-on-a-formal-group-of-height-two
\chapters@set{20261125T013500}{17}% hecke-nondegeneracy-on-the-tree
\chapters@set{20261126T030200}{18}% the-bound-is-attained
\chapters@set{20261125T009000}{19}% kronecker-limit-and-chowla-selberg
\chapters@set{20261125T002000}{20}% definite-theta-lifts
\chapters@set{20261125T015500}{21}% siegel-weil-in-the-smallest-case
\chapters@set{20261125T002100}{22}% definite-theta-lifts-and-norms
\chapters@set{20261125T016500}{23}% hecke-operators-and-theta-lifts
\chapters@set{20261125T016700}{24}% divisibility-along-the-chain
\chapters@set{20261125T004000}{25}% denominators-and-sha
\chapters@set{20261125T004100}{26}% tate-twists-and-sha
\chapters@set{20261126T020430}{27}% the-minus-part-by-gauss-sums
\chapters@set{20261125T014100}{28}% congruence-modules
\chapters@set{20261125T014500}{29}% eisenstein-congruences-and-selmer-classes
\chapters@set{20261125T014700}{30}% pseudocharacters-and-the-length-bound
\chapters@set{20261125T011000}{31}% galois-representations-black-box
\chapters@set{20261126T030300}{32}% the-data
\chapters@set{20261126T020210}{33}% bessel-poincare-series-and-their-gram-operators
\chapters@set{20261126T030400}{34}% the-seesaw-identity
\chapters@set{20261126T030500}{35}% divisibility-of-the-composite-lift
\chapters@set{20261126T030600}{36}% theorem-a-short-forms
\chapters@set{20261126T030700}{37}% the-determinant-of-a-chain
\chapters@set{20261126T030800}{38}% the-cost-of-one-lift
\chapters@set{20261126T030900}{39}% denominators-of-projections-to-lifts
\chapters@set{20261126T031000}{40}% twists-with-few-congruences
\chapters@set{20261126T031100}{41}% theorem-c-and-the-main-inequality
\providecommand\chapnumber[1]{%
  \@ifundefined{chapters@#1}%
    {\PackageError{chapters}{Sheet #1 is not listed in the index}%
      {List it in 20261125T012000--grh-workbook__guide.org\MessageBreak
       and run scripts/make-chapters.py.}}%
    {\edef\thechap{\@nameuse{chapters@#1}}}}
\makeatother
\chapnumber{<sheet ID>}
\makeatletter
\providecommand\chapters@set[2]{\@namedef{chapters@#1}{#2}}
\chapters@set{20261125T001000}{1}% zeros-to-ramp-sums
\chapters@set{20261125T001100}{2}% inert-primes-and-estimate-b
\chapters@set{20261125T005000}{3}% galois-fourier-and-averaging
\chapters@set{20261125T005500}{4}% columns-from-a-sphere
\chapters@set{20261125T000000}{5}% slope-method
\chapters@set{20261126T030100}{6}% a-twist-for-each-prime
\chapters@set{20261125T003000}{7}% the-chain-and-the-square-root
\chapters@set{20261125T010000}{8}% rank-lemma-and-zero-estimate
\chapters@set{20261125T010500}{9}% poincare-series-and-near-orthogonality
\chapters@set{20261124T223000}{10}% elliptic-curves-for-the-proof
\chapters@set{20261125T009500}{11}% fock-model-and-theta-intertwiners
\chapters@set{20261124T223100}{12}% complex-multiplication-and-inert-reduction
\chapters@set{20261125T008000}{13}% one-cm-determinant
\chapters@set{20261124T214500}{14}% supersingular-theta-divisibility
\chapters@set{20261124T214600}{15}% the-isotropic-plane
\chapters@set{20261125T013700}{16}% pairings-on-a-formal-group-of-height-two
\chapters@set{20261125T013500}{17}% hecke-nondegeneracy-on-the-tree
\chapters@set{20261126T030200}{18}% the-bound-is-attained
\chapters@set{20261125T009000}{19}% kronecker-limit-and-chowla-selberg
\chapters@set{20261125T002000}{20}% definite-theta-lifts
\chapters@set{20261125T015500}{21}% siegel-weil-in-the-smallest-case
\chapters@set{20261125T002100}{22}% definite-theta-lifts-and-norms
\chapters@set{20261125T016500}{23}% hecke-operators-and-theta-lifts
\chapters@set{20261125T016700}{24}% divisibility-along-the-chain
\chapters@set{20261125T004000}{25}% denominators-and-sha
\chapters@set{20261125T004100}{26}% tate-twists-and-sha
\chapters@set{20261126T020430}{27}% the-minus-part-by-gauss-sums
\chapters@set{20261125T014100}{28}% congruence-modules
\chapters@set{20261125T014500}{29}% eisenstein-congruences-and-selmer-classes
\chapters@set{20261125T014700}{30}% pseudocharacters-and-the-length-bound
\chapters@set{20261125T011000}{31}% galois-representations-black-box
\chapters@set{20261126T030300}{32}% the-data
\chapters@set{20261126T020210}{33}% bessel-poincare-series-and-their-gram-operators
\chapters@set{20261126T030400}{34}% the-seesaw-identity
\chapters@set{20261126T030500}{35}% divisibility-of-the-composite-lift
\chapters@set{20261126T030600}{36}% theorem-a-short-forms
\chapters@set{20261126T030700}{37}% the-determinant-of-a-chain
\chapters@set{20261126T030800}{38}% the-cost-of-one-lift
\chapters@set{20261126T030900}{39}% denominators-of-projections-to-lifts
\chapters@set{20261126T031000}{40}% twists-with-few-congruences
\chapters@set{20261126T031100}{41}% theorem-c-and-the-main-inequality
\providecommand\chapnumber[1]{%
  \@ifundefined{chapters@#1}%
    {\PackageError{chapters}{Sheet #1 is not listed in the index}%
      {List it in 20261125T012000--grh-workbook__guide.org\MessageBreak
       and run scripts/make-chapters.py.}}%
    {\edef\thechap{\@nameuse{chapters@#1}}}}
\makeatother
\chapnumber{<sheet ID>}
\makeatletter
\providecommand\chapters@set[2]{\@namedef{chapters@#1}{#2}}
\chapters@set{20261125T001000}{1}% zeros-to-ramp-sums
\chapters@set{20261125T001100}{2}% inert-primes-and-estimate-b
\chapters@set{20261125T005000}{3}% galois-fourier-and-averaging
\chapters@set{20261125T005500}{4}% columns-from-a-sphere
\chapters@set{20261125T000000}{5}% slope-method
\chapters@set{20261126T030100}{6}% a-twist-for-each-prime
\chapters@set{20261125T003000}{7}% the-chain-and-the-square-root
\chapters@set{20261125T010000}{8}% rank-lemma-and-zero-estimate
\chapters@set{20261125T010500}{9}% poincare-series-and-near-orthogonality
\chapters@set{20261124T223000}{10}% elliptic-curves-for-the-proof
\chapters@set{20261125T009500}{11}% fock-model-and-theta-intertwiners
\chapters@set{20261124T223100}{12}% complex-multiplication-and-inert-reduction
\chapters@set{20261125T008000}{13}% one-cm-determinant
\chapters@set{20261124T214500}{14}% supersingular-theta-divisibility
\chapters@set{20261124T214600}{15}% the-isotropic-plane
\chapters@set{20261125T013700}{16}% pairings-on-a-formal-group-of-height-two
\chapters@set{20261125T013500}{17}% hecke-nondegeneracy-on-the-tree
\chapters@set{20261126T030200}{18}% the-bound-is-attained
\chapters@set{20261125T009000}{19}% kronecker-limit-and-chowla-selberg
\chapters@set{20261125T002000}{20}% definite-theta-lifts
\chapters@set{20261125T015500}{21}% siegel-weil-in-the-smallest-case
\chapters@set{20261125T002100}{22}% definite-theta-lifts-and-norms
\chapters@set{20261125T016500}{23}% hecke-operators-and-theta-lifts
\chapters@set{20261125T016700}{24}% divisibility-along-the-chain
\chapters@set{20261125T004000}{25}% denominators-and-sha
\chapters@set{20261125T004100}{26}% tate-twists-and-sha
\chapters@set{20261126T020430}{27}% the-minus-part-by-gauss-sums
\chapters@set{20261125T014100}{28}% congruence-modules
\chapters@set{20261125T014500}{29}% eisenstein-congruences-and-selmer-classes
\chapters@set{20261125T014700}{30}% pseudocharacters-and-the-length-bound
\chapters@set{20261125T011000}{31}% galois-representations-black-box
\chapters@set{20261126T030300}{32}% the-data
\chapters@set{20261126T020210}{33}% bessel-poincare-series-and-their-gram-operators
\chapters@set{20261126T030400}{34}% the-seesaw-identity
\chapters@set{20261126T030500}{35}% divisibility-of-the-composite-lift
\chapters@set{20261126T030600}{36}% theorem-a-short-forms
\chapters@set{20261126T030700}{37}% the-determinant-of-a-chain
\chapters@set{20261126T030800}{38}% the-cost-of-one-lift
\chapters@set{20261126T030900}{39}% denominators-of-projections-to-lifts
\chapters@set{20261126T031000}{40}% twists-with-few-congruences
\chapters@set{20261126T031100}{41}% theorem-c-and-the-main-inequality
\providecommand\chapnumber[1]{%
  \@ifundefined{chapters@#1}%
    {\PackageError{chapters}{Sheet #1 is not listed in the index}%
      {List it in 20261125T012000--grh-workbook__guide.org\MessageBreak
       and run scripts/make-chapters.py.}}%
    {\edef\thechap{\@nameuse{chapters@#1}}}}
\makeatother

\ifcsname chapters@20261125T020400\endcsname
  \chapnumber{20261125T020400}
\else
  \renewcommand{\thechap}{?}
\fi
\renewcommand{\thesection}{\thechap.\arabic{section}}
\theoremstyle{definition}
\newtheorem{problem}{Problem}[section]
\theoremstyle{plain}
\newtheorem{theorem}{Theorem}[section]
\theoremstyle{definition}
\newtheorem{hypothesis}[theorem]{Hypothesis}
\newcommand{\Z}{\mathbb Z}
\newcommand{\Q}{\mathbb Q}
\newcommand{\F}{\mathbb F}
\newcommand{\Oc}{\mathcal O}
\newcommand{\Sha}{\mathrm{Sha}}
\newcommand{\len}{\operatorname{length}}
\newcommand{\num}{\operatorname{num}}
\setlength{\emergencystretch}{2em}

\title{Local terms, congruence defects and Selmer windows}
\author{}
\date{November 25, 2026}

\begin{document}
\maketitle
\refstepcounter{chap}\label{chap}

\paragraph{Prerequisites.}
Chapter~\ref{sm-chap} reads a local norm from a lattice exponent (Problem~\ref{sm-pr:pairing-norm}).
We use the eigenspaces of cyclotomic class groups from Chapter~\ref{ds-chap}, Theorems~\ref{ds-thm:stickelberger} and~\ref{ds-thm:HR}, and the strict Selmer groups over \(\Q\) from Chapter~\ref{ts-chap}, Theorems~\ref{ts-thm:sha-class} and~\ref{ts-thm:MW}.  Chapter~\ref{cg-chap}, Problem~\ref{cg-pr:691}, computes a projection of the weight-twelve lattice.  Theorems~\ref{ei-thm:ribet-lemma} and~\ref{pc-thm:model} construct a residual extension and bound a cotangent length, respectively.  We also recall the local correspondence and the Galois representations of regular algebraic polarized cusp forms (Theorems~\ref{gb-thm:LLC} and~\ref{gb-thm:galois}).

Why does a split prime contribute exactly half its logarithm before one sums the local exponents? And what happens when projecting onto a Hecke eigenspace introduces a denominator? The first calculation gives a main term for each lift.  To bound the denominator of the projection, the argument uses the orders of strict Selmer groups over a varying CM field.

\section{The two level primes}\label{sec:local}
Fix \(0<\varepsilon<1/52\) and a large scale \(x\).  Write \(y=x^\varepsilon\), \(b=e^y\).
There are \(n\) columns, \(f=2n\) elements in the CM type \(\Phi\) of the host field \(E\), and \(m\) lifts.  Write \(B_0:=fmby\) and let \(A=\sum_t a_t\) be the total weight of the tops.  For a finite place \(w\mid p\) of a Galois coefficient field \(L\supset E\), use the integer valuation \(v_w\), ramification index \(e_w\), residue degree \(f_w\), and weight \(\ell_w=f_w\log p/[L:\Q]\).  Thus \(\sum_{w\mid p}\ell_we_w=\log p\).

The two distinct level primes \(g_\iota=\pi_\iota\bar\pi_\iota\), \(\iota=0,1\), split in the imaginary quadratic subfield \(K\) of \(E\).  Set \(g=\log(g_0g_1)\).  Let \(u\) index a chain of lifts.  At its \(j\)-th lift the exponents \(n_{u,j}(\rho)\) have sum \(n_{u,j}\), and the polynomial model has degree \(\Delta_j\) and shift \(\kappa_j\ge0\).  Let \(D_\Gamma\) be a positive common multiple of the stabilizer orders in the lifts.  At \(w\mid g_\iota\) the local exponent is \cite[Def.~16.5]{Adr}
\[
 E_{u,j}(w)=v_{g_\iota}(D_\Gamma)-\tfrac12\Delta_j+\kappa_j+
 e_w^{-1}\sum_{\rho\in\Phi}n_{u,j}(\rho)v_w(\tau_\rho\bar\pi_\iota),
\]
where \(\tau_\rho\) extends \(\rho\) to \(L\).  Write \(T_u=\varpi_u\bar\varpi_u\) for the terminal prime, distinct from \(g_0,g_1,d\).  At \(w\mid T_u\) the exponent is \(\varepsilon^T_u\ge0\) when \(v_w(\tau_\rho\varpi_u)>0\) for every \(\rho\in\Phi\), and zero otherwise.  At the remaining primes it is \(v_p(D_\Gamma)\), with an additional error \(\varepsilon^d_j\) at \(d\) or \(\varepsilon^{\rm in}_{j,p}\) at the remaining primes in the inert loss set \(P^{\rm in}\).  Write \(E_{u,j}:=\sum_w2\ell_we_w E_{u,j}(w)\) and \(E_u=\sum_{j=1}^m E_{u,j}\).

\begin{problem}\label{pr:half-norm}
For \(K=\Q(\mathrm i)\) and \(5=(2+\mathrm i)(2-\mathrm i)\), the places above \(5\) contribute \(0\) and \(\tfrac12\log5\) to the weighted valuation of \(2-\mathrm i\) in \(K\).  In general
\(\sum_{w\mid g_\iota}\ell_wv_w(\tau_\rho\bar\pi_\iota)=\tfrac12\log g_\iota\), independently of \(L\) and \(\rho\).
\end{problem}

\begin{problem}\label{pr:level}
The contribution of \(g_\iota\) to \(E_{u,j}\) is
\((2v_{g_\iota}(D_\Gamma)-\Delta_j+2\kappa_j+n_{u,j})\log g_\iota\).
\end{problem}

\begin{theorem}[local terms, {\cite[Thm.~16.6, Cor.~16.7]{Adr}}]\label{thm:local}
Assume the local data above and (E1)--(E6) of~\cite[\S16.2]{Adr}.
These require \(\lvert n_{u,j}-\Delta_j-A\rvert\le\delta_j\) and bound
\[
 (\delta_j+2\kappa_j)g,\quad 2\varepsilon^T_u\log T_u,\quad 2\log D_\Gamma,\quad
 2|\varepsilon^d_j|\log d,\quad \sum_{p\in P^{\rm in}}2|\varepsilon^{\rm in}_{j,p}|\log p
\]
by \(C_gB_0,C_TB_0,C_\Gamma B_0,C_dB_0,C_{\rm in}B_0\), respectively.  With \(C_E\) the sum of the five constants, \(\lvert E_{u,j}-Ag\rvert\le C_E B_0\) for every \(u,j\).
\end{theorem}
\noindent\emph{Simplest incarnation:} the local exponent of a pairing in Problem~\ref{sm-pr:pairing-norm}. \emph{Status}: proved in \cite{Adr} for the stated data.
The size estimates and the bound for inert losses having an explicit formula are proved in~\cite[Lemmas~16.10--16.12]{Adr}.  The remaining orbit losses and the degrees and sizes for the constructed lifts are proved in the Lean formalization \cite{Lean}; they are \cite[Hyps.~16.13--16.14]{Adr}, not yet written out in the published expositions.

\begin{problem}\label{pr:A5}
Summation gives \(|E_u-mAg|\le e_5=C_E mB_0\), hence (A5), \(E_u\le mAg+e_5\).  With \(n,x\) fixed this allowance is proportional to \(m^2\).
\end{problem}

\section{The denominator of a projection}\label{sec:congruences}
At a place \(w\) where the local type of the lift differs from the standard type, let \(O=\Oc_{L_w}\) and \(\ell=p_w\).  For each maximal ideal \(\mathfrak m\) of the Hecke algebra acting on integral forms, let \(\Lambda_{\mathfrak m}\) be the localized lattice.  Its rational span splits into the eigenspaces \(V_G\) occurring in the lift and the complementary eigenspaces \(V_B\).  Let \(\pi_G\) project along \(V_B\).  The congruence defect of~\cite[Def.~17.1, Prop.~17.30]{Adr} is
\[
 \eta_j(w)=\frac1{e_w}\sum_{\mathfrak m}
 \len_O\frac{\pi_G\Lambda_{\mathfrak m}}{\Lambda_{\mathfrak m}\cap V_G}.
\]
Here \(\len_O(O/\ell O)=e_w\).  At a place of standard type put \(\eta_j(w)=0\).

\begin{problem}\label{pr:691}
In the classical model take \(\Lambda=M_{12}(\Z)\otimes\Z_{691}\), \(V_G=\Q_{691}\Delta\) and \(V_B=\Q_{691}E_{12}\).  The defect is \(1\), and its weighted contribution \(\ell_we_w\eta\) over \(\Q\) is \(\log691\).
\end{problem}

\begin{theorem}[Fitting inequality, {\cite[Prop.~17.27]{Adr}}]\label{thm:fitting}
Let \(O\) be a complete discrete valuation ring with fraction field \(K_w\), and \(\mathbb T\subset\prod_{z\in X}  O\) an \(O\)-algebra on a finite set \(X\).  Let \(\delta_{z_0}\) be the least valuation of a nonzero entry at \(z_0\) of an element supported there, or infinity if there is none.  Let \(P\) be an \(a'\)-by-\(s\) matrix over an \(O\)-algebra \(R'\), \(a'\ge s\), with representations \(\rho_z\) over \(O\).
Assume (F1)--(F4).  For every \(s\)-by-\(s\) matrix \(Q\) over \(R'\), the function \(z\mapsto\det\rho_z(Q)\) belongs to \(\mathbb T\).  Each \(\rho_z(P)\) for \(z\ne z_0\) has nonzero rational kernel.  The representation \(\rho_{z_0}\) is block upper triangular, with diagonal blocks \(\rho_i\) of sizes \(d_i\) whose joint reduction maps onto the product of the full matrix algebras.  The residue field has size \(>s\sum_i d_i\).  Then
\[
 \delta_{z_0}\le\sum_i\len_O\ker\bigl(\rho_i(P):(K_w/O)^{sd_i}\longrightarrow(K_w/O)^{a'd_i}\bigr).
\]
\end{theorem}
\noindent\emph{Simplest incarnation:} the two eigenvalues in the example following Theorem~\ref{pc-thm:model}. \emph{Status}: proved in \cite{Adr}.

\begin{problem}\label{pr:fitting}
Let \(O=\Z_{691}\) and \(R'=\mathbb T=\{(v_1,v_2)\in O^2:v_1-v_2\in691O\}\), with its two coordinate representations.  For the one-entry matrix \(P=((691,0))\) and \(z_0\) the first coordinate, both sides of Theorem~\ref{thm:fitting} are \(1\).  The graded cotangent module in the weight-twelve example following Theorem~\ref{pc-thm:model} thus has the same length.
\end{problem}

\section{Strict Selmer groups and the two signs}\label{sec:selmer}
For a number field \(F\) use \(W_i=(\Q_\ell/\Z_\ell)(-i)\) with cyclotomic action to the power \(-i\).  The group \(\Sha_\ell(F,i)\) consists of continuous \(H^1(F,W_i)\) classes whose restrictions to inertia at every finite place and to the decomposition group at every place over \(\ell\) vanish.  For a finite \(\ell\)-group \(U\) put \(\lambda(U)=\len_{\Z_\ell}(U)\log\ell=\log|U|\).  The sum used in the window \cite[Def.~17.2, Hyp.~17.11]{Adr} is
\[
 V(F,t)=\sum_{\ell}\sum_{1\le i<t}\lambda(\Sha_\ell(F,i)).
\]

\begin{theorem}[restriction and class groups, {\cite[Props.~17.7--17.8]{Adr}}]\label{thm:class}
Let \(\ell\) be odd, \(i\ge1\), and \((\ell-1)\nmid i\).  For \(E\subset\Q(\mu_N)\), \(\ell\nmid N\), restriction embeds \(\Sha_\ell(\Q,i)\) into \(\Sha_\ell(E,i)\) with Galois-invariant image.  Let \(A_\ell\) be the \(\ell\)-part of the class group of \(\Q(\mu_\ell)\), and \(\omega\) the Teichm\"uller character.  The dimension of \(\Sha_\ell(\Q,i)[\ell]\) is the \(\ell\)-rank of \(A_\ell^{(\omega^{-i})}\).
\end{theorem}
\noindent\emph{Simplest incarnation:}  Theorem~\ref{ts-thm:sha-class} identifies the two sides by duality. \emph{Status}: classical.

\begin{problem}\label{pr:character}
In the Artin description of Theorem~\ref{ts-thm:sha-class}, an equivariant homomorphism \(\phi:A_\ell\to W_i[\ell]\) satisfies \(\phi(\sigma z)=\omega(\sigma)^{-i}\phi(z)\).  After choosing a basis of \(W_i[\ell]\), its possible values give precisely \(\operatorname{Hom}(A_\ell^{(\omega^{-i})},\F_\ell)\).  For even \(i\) the character is even.
\end{problem}

\begin{problem}\label{pr:odd}
For positive odd \(i\), Theorem~\ref{ts-thm:MW} gives
\(\sum_{\ell\ {\rm odd}}\lambda(\Sha_\ell(\Q,i))\le\log|\num(B_{i+1}/(i+1))|\le(i+1)\log(i+1)+3\).
For \(t=\lfloor x^{401}\rfloor\) the sum over odd \(1\le i<t\) is \(O(t^2\log(2t))=o(\sqrt b)\).
\end{problem}

For comparison, statement (V) of~\cite[\S17.4]{Adr} asserts the existence of constants \(C,P_0\) with the following property.  For every even \(i\ge2\) and finite set \(S\) of primes \(\ell\ge5\) with \((\ell-1)\nmid2i\), the groups below are finite and
\[
 \sum_{\ell\in S}\lambda(\Sha_\ell(\Q,i))\le C(i+1)^{P_0}.
\]

\begin{problem}\label{pr:vanishing}
Assuming (V), \(A_\ell^{(\omega^{-i})}=0\) once \(\log\ell>C(i+1)^{P_0}\) for these \(i,\ell\) \cite[Cor.~17.15]{Adr}.  The threshold is polynomial in \(i\) for the logarithm of the prime.
\end{problem}

Kurihara proved vanishing at \(i=2\).  For \(i+1=\nu\) odd, Soul\'e's sufficient threshold for \(\log\ell\) is \(\nu^{224\nu^4}\).  The estimates of de la Torre Dur\'an give \(\exp(O(\nu^2\log\nu))\) in the same comparison \cite[\S17.4, Rem.~17.16]{Adr}, \cite{Kur,So2,DT}.  The odd-twist estimate in Problem~\ref{pr:odd} concerns \(\Q\).  The window requires both signs over \(E\), and includes \(\ell=2\).

\begin{problem}\label{pr:ribet}
In the nonsplit residual extension of Theorem~\ref{ei-thm:ribet-lemma}, write the diagonal characters as \(\chi_1,\chi_2\) and the upper entry as \(b_{12}(\sigma)\).  Then \(b_{12}(\sigma)/\chi_2(\sigma)\) represents a nonzero class with coefficients in \(\chi_1\chi_2^{-1}\).  For the absolute Galois group of \(F\), if this coefficient line is \(W_i[\ell]\), the image of the class in \(H^1(F,W_i)\) belongs to \(\Sha_\ell(F,i)\) precisely when that image satisfies the two local conditions.
\end{problem}

In (G1) of~\cite[\S17.4]{Adr} a Selmer class gives a tangent vector at the Eisenstein point.  Passing from such vectors to an automorphic count also requires (G3) and (G4).  The comparison there with Allen--Newton--Thorne uses the distinction between individually conjugate self-dual constituents and a pair exchanged by conjugate duality.  Newton--Thorne's adjoint vanishing requires an enormous image.  Kato's vanishing applies at the reducible modular points.

\section{The window and the congruence bound}\label{sec:window}
Fix a primitive character \(\psi_0\) of conductor \(M_\psi\) and order \(o\).  For an integer \(h\ge0\) set \(N=\operatorname{lcm}(2^{h+2},M_\psi,q,d)\).  Write \(E(\psi_0,h,q,d)\) for the subfield of \(\Q(\mu_N)\) fixed by classes \(a\equiv\pm1\bmod 2^{h+2}\) with \(\psi_0(a)=1\) and \((\frac aq)=(\frac ad)=1\).  A window pair at \(Y_0\) consists of primes \(q\equiv3\bmod8,\ d\equiv7\bmod16\) with \(z=Y_0^{\varepsilon/5}\),
\[
 z/2<\log q\le z,\qquad
 \tfrac12z\log^2Y_0<\log d\le z\log^2Y_0,
\]
such that \(q\) and every odd prime \(p\le z\) divide \(d+1\) \cite[Defs.~17.9--17.10]{Adr}.

\begin{hypothesis}[Selmer window, {\cite[Hyp.~17.11]{Adr}}]\label{hyp:window}
Write \(\alpha=2+8\varepsilon\).  For each \(\psi_0,\varepsilon\) and integer-valued function \(h\ge0\) with \(x^\alpha\le2^{h(x)}\le2x^\alpha\) for \(x\ge1\), there are \(J_1\ge1,X_0\) such that for every \(X\ge X_0\) and every window pair at \(Y_0=X^{1/J_1}\) there is \(W\in[Y_0,\sqrt X]\) with the following property.  For every \(W\le x\le W^2\), with \(E=E(\psi_0,h(x),q,d)\) and \(t=\lfloor x^{401}\rfloor\), all the groups defining \(V(E,t)\) are finite and \(V(E,t)\le\sqrt b\).
\end{hypothesis}
\noindent\emph{Simplest incarnation:} the odd twists over \(\Q\) in Problem~\ref{ts-pr:one-number}. \emph{Status}: proved in the Lean formalization; it is \cite[Hyp.~17.11]{Adr}, not yet written out in the published expositions.

\begin{theorem}[the two signs, {\cite[Thm.~17.12]{Adr}}]\label{thm:signs}
For odd \(\ell\) let \(P_\pm\) sum \(\lambda\) of the \(c=\pm(-1)^i\) eigenspaces over these primes and \(1\le i<t\), with \(c\) complex conjugation.  With \(P_2:=\sum_{1\le i<t}\len_{\Z_2}\Sha_2(E,i)\), the two idempotents \((1\pm c)/2\) at odd primes give \(V(E,t)=(\log2)P_2+P_++P_-\).  Write \(F_E=\operatorname{cond}(E)\) and \(d_E=[E:\Q]\).  The following two estimates, including finiteness, imply Hypothesis~\ref{hyp:window}.
\begin{itemize}
\item[(I)] Fix \(\psi_0,\varepsilon\), and functions \(h,t\) with \(h(x)\ge2,\ 2^{h(x)}\le x^4,\ 2\le t(x)\le x^{401}\) for \(x\ge2\).  There are \(J_1,C,P_0',X_1>0\) such that for \(X\ge X_1\) and primes \(q,d\) in the indicated residue classes, prime to \(M_\psi\), with
\[
 q,d>\max(P_0',2X^{401}+1),\qquad \log(qd)\le X^{\varepsilon/(4J_1)},
\]
there is \(W\in[X^{1/J_1},\sqrt X]\) such that for \(W\le x\le W^2\),
\[
 P_2\le e^{x^{4\varepsilon/5}},\qquad P_+\le(F_Et)^C+C e^{x^{4\varepsilon/5}}.
\]
Here \(E=E(\psi_0,h(x),q,d)\) and \(t=t(x)\).
\item[(II)] For all such fields and \(t\ge2\), with an absolute constant \(C_-\),
\[
 P_-\le 3d_Et^2\log(tF_E)+C_-\sum_{\substack{\ell\mid o\\\ell\ {\rm odd}}}
          \ell d_E^2t^2\log^2(t\ell F_E).
\]
\end{itemize}
\end{theorem}
\noindent\emph{Simplest incarnation:} the minus eigenspaces of Theorem~\ref{ds-thm:HR}. \emph{Status}: proved in \cite{Adr} as a reduction.  For the restriction of a rational class, \(c\) acts trivially, so even twists occur in \(P_+\) and odd twists in \(P_-\).  The estimates (I) and (II) are inputs of this reduction, and are not proved in the published expositions.  The minus estimate is of abelian main-conjecture type.  The generic plus estimate and (V) both use the character estimate (CorGK), for \(K=\Q(\sqrt{-3})\) in the deduction of (V) \cite[Props.~17.14, 17.19]{Adr}.

\begin{problem}\label{pr:conductor}
For \(h\) as in Hypothesis~\ref{hyp:window}, \(x\ge Y_0\) and \(t=\lfloor x^{401}\rfloor\), the inequalities for the window pair give \(\log(F_Et)\le 2x^{\varepsilon/4}\) for all sufficiently large \(x\).  Consequently each fixed power of \(F_Et\) and \(e^{x^{4\varepsilon/5}}\) is \(o(\sqrt b)\).
\end{problem}

\begin{problem}\label{pr:window-consequence}
Fix \(\psi_0,\varepsilon,h\) as in Hypothesis~\ref{hyp:window} and assume it.  Suppose window pairs exist at \(X^{1/J_1}\) for unbounded \(X\).  Then for arbitrarily large \(x\) the sum of \(\lambda(\Sha_\ell(\Q,i))\) over odd \(\ell\nmid M_\psi q_xd_x\) and \(1\le i<\lfloor x^{401}\rfloor\) with \((\ell-1)\nmid i\) is at most \(\sqrt b\).  The two excluded window primes may vary with \(x\) \cite[Prop.~17.22]{Adr}.
\end{problem}

What uniform estimate would an order formula for these groups have to supply?  Burns--Greither and Huber--Kings treat equivariant special values for abelian fields and Dirichlet characters \cite{BuG,HK}.

\begin{problem}\label{pr:order-formula}
Suppose that for a given \(E,t\) every \(\Sha_\ell(E,i)\), \(1\le i<t\), has order \(\ell^{v_\ell(a_{\ell,i})}\), with nonzero \(a_{\ell,i}\in\Z_\ell\).  Then the numerical condition in Hypothesis~\ref{hyp:window} is exactly
\(\sum_{\ell}\sum_{1\le i<t}v_\ell(a_{\ell,i})\log\ell\le\sqrt b\)
at each scale in its window.  Finiteness of the individual groups alone gives no such upper bound.
\end{problem}

\begin{theorem}[rescaled window, {\cite[Props.~18.3--18.4]{Adr}}]\label{thm:W3}
Assume Hypothesis~\ref{hyp:window}.  In the construction \(n=c_*2^{h(x)}\in[x^\alpha,2x^\alpha]\) with \(c_*\mid o\) eventually constant.  There is \(J\ge2\) and there are arbitrarily large \(Y\) with base \(Y_0\le Y,\ Y^2\le Y_0^J\) and a window pair at \(Y_0\) such that throughout \(Y\le x\le Y^{3/2}\),
\[
 V(E,\lfloor x^{400}\rfloor)\le\sqrt b,\qquad E=E(\psi_0,h(x),q,d).
\]
All groups in this sum are finite.  With a sufficiently large base the same bound holds for all integer twists \(1\le i<x^{400}\), as in (W3).
\end{theorem}
\noindent\emph{Simplest incarnation:} the single-field sum in Problem~\ref{ts-pr:one-number}. \emph{Status}: proved in \cite{Adr} from Hypothesis~\ref{hyp:window}, which is proved in the Lean formalization and is not yet written out in the published expositions.  The prime selections in~\cite[Thms.~5.10, 5.12, Prop.~18.4]{Adr} supply pairs at arbitrarily large bases.

\begin{problem}\label{pr:rescale}
Write \(\theta=c_*^{1/\alpha}\) and apply Hypothesis~\ref{hyp:window} with exponent \(h(\theta s)\) for sufficiently large \(s\).  For \(Y=\theta W\) and \(Y_0\ge\max(3,2\theta^{400})\), one has \(Y^2\le Y_0^{J_1+1}\).  If \(Y\le x\le Y^{3/2}\), then \(s=x/\theta\in[W,W^2]\) and \(2x^{400}\le s^{401}\).  Thus every integer \(i<x^{400}\) lies below \(\lfloor s^{401}\rfloor\).
\end{problem}

Let \(k\) be the rank parameter for the lifts, \(T_2\) the distinguished set of columns, and \(N_c=|T_2|kby^2/n^2\).  Let \(t_r\) count the copies of the auxiliary space added at each lift.  The ranks in the construction give \(t_r\le x^{400}-1\) \cite[Lemma~18.2]{Adr}.  Write \(d_j^0\) for the rank of lifted forms of standard type at step \(j\).

\begin{hypothesis}[congruences at the twisted places, {\cite[Hyp.~17.24]{Adr}}]\label{hyp:congruence}
There are constants \(C_N,x_0\) such that for every admissible datum at \(x\ge x_0\) and every \(j\), if the groups in \(V(E,t_r)\) are finite and the sum converges, then
\[
 \sum_{w\ {\rm twisted}}\ell_we_w\eta_j(w)
 \le d_j^0\{C_N(yB_0+(\log x)^2N_c)+V(E,t_r)\}.
\]
\end{hypothesis}
\noindent\emph{Simplest incarnation:} the projection denominator in Problem~\ref{cg-pr:691}.  Theorem~\ref{cg-thm:hida} gives the classical cuspidal comparison. \emph{Status}: proved in the Lean formalization; it is \cite[Hyp.~17.24]{Adr}, not yet written out in the published expositions, and \cite{Adr} outlines its proof as follows.
The lattice quotient is a sum of projection defects, bounded using the congruence exponent of each eigensystem with its multiplicity \cite[Props.~17.25, 17.30]{Adr}.  After removing eigensystems outside the chosen local types, the per-component inequalities and their weighted sum yield the congruence bound, under the multiplicity and size bounds in~\cite[Prop.~17.34, Hyp.~17.36]{Adr}.  The Galois representations of Theorem~\ref{gb-thm:galois}, Poitou--Tate duality and Berger's potential semistability theorem enter the comparison in~\cite[\S17.5]{Adr}.  The estimate for lifting crystalline classes with Wach modules, the other local estimates and the integral comparison with theta series are not yet written out in the published expositions \cite[\S18.7]{Adr}.

\begin{problem}\label{pr:budget}
Assume Hypothesis~\ref{hyp:congruence}, (W3), \(d_j^0\le D_0\) and \((\log x)^2\le y\).  Then \(\sum_w\ell_we_w\eta_j(w)\le (C_N+1)D_0y(B_0+N_c)\) \cite[Thm.~17.37]{Adr}.  Dividing by the positive rank bound \(D_0\) and summing \(m\) lifts puts \(mN_c\), with the common factor \(y\), in the height allowance.  The formula for \(e_5\) in Problem~\ref{pr:A5} uses only \(B_0\).
\end{problem}

\begin{theorem}[cost of the congruence defects, {\cite[Thm.~17.37]{Adr}}]\label{thm:congruence}
On an admissible window satisfying (W3) and Hypothesis~\ref{hyp:congruence}, if \(d_j^0\le D_0\) for all \(j\) and \((\log x)^2\le y\), each lift has weighted congruence defect at most \((C_N+1)D_0y(B_0+N_c)\).
\end{theorem}
\noindent\emph{Simplest incarnation:} the single summand at \(691\) in Problem~\ref{cg-pr:691}. \emph{Status}: proved in \cite{Adr} under the indicated hypotheses.

\section*{Hints}
\noindent\ref{pr:half-norm}.  Use the norm from \(L\) to \(\Q\) of an element of \(K\).\\
\ref{pr:level}.  Apply the two sums over \(w\mid g_\iota\) to the two parts of the local exponent.\\
\ref{pr:A5}.  Keep \(B_0\) in the sum until its last occurrence.\\
\ref{pr:691}.  Localize the two lattices in Problem~\ref{cg-pr:691}.\\
\ref{pr:fitting}.  Find the elements supported on the first coordinate.  Use \((K_w/O)[691]\) on the other side.\\
\ref{pr:character}.  Apply equivariance to each character summand of \(A_\ell/\ell A_\ell\).\\
\ref{pr:odd}.  Use Problems~\ref{ts-pr:one-number} and~\ref{ts-pr:bernoulli-size}, then compare logarithms.\\
\ref{pr:vanishing}.  Apply (V) to a singleton prime.\\
\ref{pr:ribet}.  Multiply two upper triangular matrices and then change the basis of the extension.\\
\ref{pr:conductor}.  Use \(F_E\le N\le 2^{h+2}M_\psi qd\) and \(\varepsilon/4-\varepsilon/5>0\).\\
\ref{pr:window-consequence}.  Take \(x=W\) and use Theorem~\ref{thm:class} prime by prime.\\
\ref{pr:order-formula}.  Test the size condition on a single cyclic group.\\
\ref{pr:rescale}.  Compare \(\theta\) and \(\theta^{400}\) with the lower endpoint of the window.\\
\ref{pr:budget}.  Compare \(\sqrt b\) with \(B_0\) before summing over \(j\).

\medskip
\paragraph{Deliverable.}\label{par:deliverable}
For the chain record of~\cite[Def.~18.9, Prop.~18.14]{Adr}, Theorem~\ref{thm:local} gives (K5) under the degree and orbit hypotheses, and Problem~\ref{pr:A5} gives (A5).  Theorem~\ref{thm:congruence} supplies the congruence contribution \(mN_c\) in the error of (K4), hence of (A4), under the component and rank bounds.  Theorem~\ref{thm:W3} supplies the independent window input (W3) from Hypothesis~\ref{hyp:window}, which is proved in the Lean formalization.

\begin{thebibliography}{Adr}
\bibitem[Adr]{Adr}  J.~Adriazola, \emph{The Alp\"oge proof of the Riemann hypothesis}, manuscript, revised 24 November 2026, \url{https://github.com/jimjamzola/RH-Alpoge-Reconstruction}.
\bibitem[BuG]{BuG}  D.~Burns and C.~Greither, \emph{On the equivariant Tamagawa number conjecture for Tate motives}, Invent.  Math. \textbf{153} (2003), 303--359.
\bibitem[DT]{DT}  A.~de la Torre Dur\'an, \emph{Perfect forms and torsion bounds in algebraic \(K\)-theory over totally real and CM fields}, \href{https://arxiv.org/abs/2611.06300}{arXiv:2611.06300}.
\bibitem[HK]{HK}  A.~Huber and G.~Kings, \emph{Bloch--Kato conjecture and main conjecture of Iwasawa theory for Dirichlet characters}, \href{https://arxiv.org/abs/math/0101071}{arXiv:math/0101071}.
\bibitem[Kur]{Kur}  M.~Kurihara, \emph{Some remarks on conjectures about cyclotomic fields and \(K\)-groups of \(\Z\)}, Compositio Math. \textbf{81} (1992), 223--236.
\bibitem[Lean]{Lean} The Lean community, \emph{riemann-hypothesis}, a Lean formalization of Alp\"oge's argument, \url{https://github.com/leanprover-community/riemann-hypothesis}.
\bibitem[So2]{So2}  C.~Soul\'e, \emph{Perfect forms and the Vandiver conjecture}, J.~Reine Angew.  Math. \textbf{517} (1999), 209--221.
\end{thebibliography}
\end{document}
